\documentclass[a4paper]{article}

%------------------------------------------------
% الحزم
%------------------------------------------------

\usepackage{enumitem}
\usepackage[margin=2.2cm]{geometry}
\usepackage{amsmath,amssymb}
\usepackage{graphicx}
\usepackage{array}
\usepackage{fancyhdr}
\usepackage[table]{xcolor}
\usepackage{tikz}
\usetikzlibrary{positioning,arrows.meta,shapes.geometric,calc}
\usepackage[skins]{tcolorbox}
\usepackage{polyglossia}

%------------------------------------------------
% اللغات والخطوط
%------------------------------------------------

\setdefaultlanguage[
    calendar=gregorian,
    numerals=maghrib
]{arabic}

\setotherlanguage{english}

% إذا لم تكن الخطوط مثبتة في الجهاز، تُؤخذ من المجلد fonts/
\IfFontExistsTF{Amiri}{%
    \newfontfamily\arabicfont[Script=Arabic, Scale=1.15]{Amiri}
}{%
    \newfontfamily\arabicfont[
        Script=Arabic, Scale=1.15,
        Path=fonts/, Extension=.ttf,
        UprightFont=*-Regular, BoldFont=*-Bold,
        ItalicFont=*-Italic, BoldItalicFont=*-BoldItalic
    ]{Amiri}
}

\IfFontExistsTF{AlBattar}{%
    \newfontfamily\arabicfontsf[Script=Arabic, Scale=1.2]{AlBattar}
    % خط العناوين (يعمل حتى بعد استعمال \textenglish)
    \newfontfamily\khatt[Script=Arabic, Scale=1.2]{AlBattar}
}{%
    \newfontfamily\arabicfontsf[
        Script=Arabic, Scale=1.2, Path=fonts/,
        BoldFont=ae_AlBattar.ttf
    ]{ae_AlBattar.ttf}
    % خط العناوين (يعمل حتى بعد استعمال \textenglish)
    \newfontfamily\khatt[
        Script=Arabic, Scale=1.2, Path=fonts/,
        BoldFont=ae_AlBattar.ttf
    ]{ae_AlBattar.ttf}
}

%------------------------------------------------
% الألوان: لون رئيسي هادئ + لون تنبيه فقط
%------------------------------------------------

\definecolor{asasi}{HTML}{3D6A8E}   % اللون الرئيسي: للرسوم والمخططات
\definecolor{tanbih}{HTML}{C0603A}  % لون التنبيه: للتحذيرات والحالات السلبية
\definecolor{ramadi}{HTML}{5F6B73}  % رمادي للنصوص الثانوية والأسهم
\definecolor{dakin}{HTML}{263238}   % لون العناوين

% أسماء قديمة تُحوَّل إلى اللونين فقط
\colorlet{azraq}{asasi}
\colorlet{fayruzi}{asasi}
\colorlet{banafsaji}{asasi}
\colorlet{akhdar}{asasi}
\colorlet{ahmar}{tanbih}
\colorlet{burtuqali}{asasi}

%------------------------------------------------
% الصناديق
%------------------------------------------------

\tcbset{
    asl/.style={
        enhanced,
        arc=2mm,
        boxrule=0.6pt,
        colback=white,
        colframe=black!18,
        colbacktitle=black!5,
        coltitle=dakin,
        titlerule=0pt,
        left=8pt, right=8pt, top=5pt, bottom=5pt,
        fonttitle=\khatt\bfseries,
        before skip=8pt, after skip=8pt
    }
}

% أهداف الدرس
\newtcolorbox{ahdaf}{asl, title={أهداف الدرس}}

% تعريف (خط رئيسي رفيع على الجانب)
\newtcolorbox{taarif}[1][تعريف]{asl, title={#1},
    borderline east={2pt}{0pt}{asasi}}

% مثال
\newtcolorbox{mithal}[1][مثال]{asl, title={#1}}

% تنبيه (اللون الوحيد البارز في النص)
\newtcolorbox{intabih}[1][انتبه!]{asl, colback=tanbih!5, colframe=tanbih,
    colbacktitle=tanbih, coltitle=white, title={#1}}

% ركن أستاذ اللغة العربية
\newtcolorbox{arabiyya}[1][في قسم اللغة العربية]{asl, title={#1}}

% نص أمر يكتب للأداة
\newtcolorbox{amr}[1][ramadi]{enhanced, arc=1mm, boxrule=0pt,
    colback=black!4, borderline east={2.5pt}{0pt}{#1},
    left=8pt, right=10pt, top=4pt, bottom=4pt,
    before skip=6pt, after skip=6pt}

% الخلاصة
\newtcolorbox{khulasa}{asl, title={خلاصة الدرس}}

% عنوان القسم
\newcounter{qism}
\newcommand{\qism}[1]{%
    \stepcounter{qism}%
    \par\bigskip
    {\large\khatt\bfseries{\color{asasi}\arabic{qism}.}\ {\color{dakin}#1}}\par
    \vspace{-2pt}{\color{black!20}\rule{\linewidth}{0.6pt}}\par\nopagebreak[4]
    \smallskip}

% عنوان فرعي
\newcommand{\far}[1]{\par\smallskip
    {\color{dakin}\khatt\bfseries {\color{asasi}$\blacktriangleleft$}\ #1}\par\nopagebreak[4]}

% كلمة مهمة (خط عريض فقط)
\newcommand{\kalima}[2][]{\textbf{#2}}

%------------------------------------------------
% أنماط الرسوم
%------------------------------------------------

\newcommand{\nodecolor}{}

\tikzset{
    % كتابة النص العربي من اليمين إلى اليسار داخل العقد
    % (لون النص يُعطى بالمفتاح lawn=... وليس text=...)
    every node/.append style={
        execute at begin node={\nodecolor\everypar{\beginR}\beginR},
        execute at end node=\endR},
    lawn/.code={\def\nodecolor{\color{#1}}},
    khatwa/.style={draw=#1, fill=#1!12, very thick, rounded corners=6pt,
        minimum height=1.2cm, text width=2.3cm, align=center,
        font=\khatt\bfseries},
    khatwa/.default=azraq,
    sahm/.style={-{Stealth[length=3mm]}, very thick, color=ramadi},
    raqm/.style={circle, fill=#1, lawn=white, inner sep=2pt,
        font=\small\bfseries},
}

%------------------------------------------------
% إعداد الصفحة
%------------------------------------------------

\renewcommand{\footrulewidth}{1pt}
\renewcommand{\headrulewidth}{1pt}

\rhead{\khatt المدرسة العليا للأساتذة ببوزريعة}
\lhead{\khatt ملحقة بومرداس}

\cfoot{\thepage}
\lfoot{\khatt الأستاذ: د. بن عبد الرزاق يوسف}

\parindent 0pt
\setlength{\headsep}{15pt}
\setlength{\parskip}{4pt}

\pagestyle{fancy}

\setlist[enumerate,1]{label=\arabic*), itemsep=2pt, topsep=3pt}
\setlist[itemize]{itemsep=2pt, topsep=3pt}

% الجداول
\arrayrulecolor{black!22}
\renewcommand{\arraystretch}{1.35}

% رأس جدول ملون
\newcommand{\ras}[1]{\cellcolor{black!7}\color{dakin}\bfseries #1}

% عنوان الدرس
\newcommand{\unwan}[3]{%
    \begin{center}
        {\color{asasi}\rule{\linewidth}{1.2pt}}\par\medskip
        {\khatt\color{ramadi} #1}\par\smallskip
        {\Large\khatt\bfseries\color{dakin} #2}\par\smallskip
        {\khatt\color{ramadi} #3}\par\medskip
        {\color{asasi}\rule{\linewidth}{1.2pt}}
    \end{center}}

%================================================
% بداية الوثيقة
%================================================

\begin{document}

%------------------------------------------------
% العنوان
%------------------------------------------------

\unwan{درس الفصل الثالث}
{التفاعل الفعّال مع الذكاء الاصطناعي}
{مبادئ هندسة الأوامر \textenglish{(Prompt Engineering)}}

\begin{ahdaf}
في نهاية هذا الدرس، تستطيع أن:
\begin{itemize}
    \item تعرّف الأمر وهندسة الأوامر.
    \item تذكر العناصر الخمسة للأمر الجيد.
    \item تحوّل أمرًا عامًا إلى أمر دقيق خطوة بخطوة.
    \item تستخرج عناصر أي أمر مكتوب.
\end{itemize}
\end{ahdaf}

%================================================
\qism{وضعية للتفكير}
%================================================

كتب طالبان طلبًا في الموضوع نفسه:

\begin{amr}[ahmar]
\textbf{الطالب الأول:} أعطني خطة للمراجعة.
\end{amr}

\begin{amr}[akhdar]
\textbf{الطالب الثاني:} تصرّف كمساعد أكاديمي.
أنشئ خطة مراجعة لمادة النحو مدتها 7 أيام، بمعدل ساعتين يوميًا.
حدّد موضوع كل يوم، وأضف تمرينًا قصيرًا، ونظّم الخطة في جدول.
\end{amr}

الطالب الثاني سيحصل على إجابة أفضل بكثير.
لماذا؟ لأنه \kalima{أعطى الأداة كل ما تحتاجه} لتفهم طلبه.

%================================================
\qism{ما هندسة الأوامر؟}
%================================================

\begin{taarif}
\textbf{الأمر} \textenglish{(Prompt)}: النص الذي نكتبه للأداة لنطلب منها شيئًا.

\textbf{هندسة الأوامر}: مهارة كتابة أوامر واضحة ومنظمة،
للحصول على أفضل إجابة ممكنة.
\end{taarif}

الأداة لا تعرف ما في ذهنك.
هي تعرف فقط ما تكتبه لها، لذلك يجب أن يكون الأمر واضحًا وكاملًا.

%================================================
\qism{العناصر الخمسة للأمر الجيد}
%================================================

\begin{center}
\begin{tikzpicture}[
    unsur/.style={draw=#1, fill=#1!12, very thick, circle,
        minimum size=2.4cm, align=center, font=\khatt\bfseries}]
    \node[draw=azraq, fill=azraq!20, ultra thick, rounded corners=8pt,
        minimum width=3cm, minimum height=1.3cm,
        font=\khatt\bfseries\large] (m) {الأمر الجيد};
    \node[unsur=azraq]     (r) at (90:3.3)   {الدور};
    \node[unsur=akhdar]    (t) at (18:3.6)   {المهمة};
    \node[unsur=burtuqali] (s) at (-54:3.4)  {السياق};
    \node[unsur=asasi]     (q) at (-126:3.4) {القيود};
    \node[unsur=banafsaji] (f) at (162:3.6)  {التنسيق};
    \foreach \x in {r,t,s,q,f} \draw[very thick, ramadi] (m) -- (\x);
\end{tikzpicture}
\end{center}

\begin{center}
\begin{tabular}{>{\centering\arraybackslash}m{2.2cm}
                >{\raggedleft\arraybackslash}m{6.2cm}
                >{\raggedleft\arraybackslash}m{5.6cm}}
\hline
\ras{العنصر} & \ras{السؤال الذي يجيب عنه} & \ras{مثال} \\
\hline
\kalima{الدور}            & من تكون الأداة؟ أي خبرة تعتمد؟ & تصرّف كأستاذ لغة عربية. \\
\hline
\kalima[akhdar]{المهمة}    & ماذا نريد أن تفعل الأداة؟ & اشرح درس الحال. \\
\hline
\kalima[burtuqali]{السياق} & ما المعلومات التي تساعدها على الفهم؟ & التلاميذ في السنة الثانية متوسط. \\
\hline
\kalima[ahmar]{القيود}     & ما الشروط التي يجب احترامها؟ & في 150 كلمة، بلغة بسيطة. \\
\hline
\kalima[banafsaji]{التنسيق} & كيف نريد شكل الإجابة؟ & تعريف، ثم أمثلة، ثم خلاصة. \\
\hline
\end{tabular}
\end{center}

%================================================
\qism{بناء الأمر خطوة بخطوة}
%================================================

نبدأ بأمر بسيط، ثم نضيف عنصرًا في كل مرة.
لاحظ كيف يصبح الأمر أدق في كل درجة:

\begin{center}
\begin{tikzpicture}[
    daraja/.style={draw=#1, fill=#1!12, very thick, rounded corners=4pt,
        minimum height=0.9cm, align=center, font=\khatt\bfseries}]
    \node[daraja=ahmar,     minimum width=3cm, anchor=south east] at (0,0)      {أمر عام};
    \node[daraja=burtuqali, minimum width=3cm, anchor=south east] at (-3.2,0.9) {+ السياق};
    \node[daraja=fayruzi,   minimum width=3cm, anchor=south east] at (-6.4,1.8) {+ المهمة والقيود};
    \node[daraja=akhdar,    minimum width=3cm, anchor=south east] at (-9.6,2.7) {+ التنسيق};
    \draw[sahm, azraq] (-0.3,1.3) -- (-11.5,4.3)
        node[pos=0.4, above, sloped, font=\khatt\bfseries, lawn=azraq] {جودة أفضل};
\end{tikzpicture}
\end{center}

\far{الدرجة 1: أمر عام}
\begin{amr}[ahmar]
اشرح الحال.
\end{amr}

\far{الدرجة 2: نضيف السياق}
\begin{amr}[ramadi]
أنا أستاذ متربص. تلاميذي في السنة الثانية متوسط.
اشرح لي درس الحال.
\end{amr}

\far{الدرجة 3: نضيف المهمة والقيود}
\begin{amr}[asasi]
أنا أستاذ متربص. تلاميذي في السنة الثانية متوسط.
اشرح درس الحال في 150 كلمة، بلغة بسيطة، مع مثالين من الحياة اليومية.
\end{amr}

\far{الدرجة 4: نضيف التنسيق}
\begin{amr}[akhdar]
أنا أستاذ متربص. تلاميذي في السنة الثانية متوسط.
اشرح درس الحال في 150 كلمة، بلغة بسيطة، مع مثالين من الحياة اليومية.
نظّم الإجابة في ثلاثة أجزاء: تعريف، أمثلة، خلاصة.
\end{amr}

%================================================
\qism{كيف نقرأ أمرًا ونستخرج عناصره؟}
%================================================

نقسم الأمر إلى أجزاء، ونحدد العنصر الذي يمثله كل جزء:

\begin{amr}[asasi]
تصرّف كأستاذ جامعي في البلاغة.
اشرح مفهوم التشبيه لطالب في السنة الأولى، في أقل من 200 كلمة،
مع جدول يبيّن أركان التشبيه.
\end{amr}

\begin{center}
\begin{tabular}{>{\raggedleft\arraybackslash}m{8cm}
                >{\centering\arraybackslash}m{3cm}}
\hline
\ras{الجزء من الأمر} & \ras{العنصر} \\
\hline
تصرّف كأستاذ جامعي في البلاغة.  & الدور \\ \hline
اشرح مفهوم التشبيه              & المهمة \\ \hline
لطالب في السنة الأولى           & السياق \\ \hline
في أقل من 200 كلمة              & القيود \\ \hline
مع جدول يبيّن أركان التشبيه     & التنسيق \\ \hline
\end{tabular}
\end{center}

%================================================
\qism{أخطاء شائعة}
%================================================

\begin{intabih}
\begin{itemize}
    \item \textbf{أمر قصير جدًا:} «اشرح البلاغة» لا يحدد المستوى ولا الهدف.
    \item \textbf{أمر غامض:} «اكتب شيئًا جميلًا» لا يحدد ما المطلوب.
    \item \textbf{عدة مهام مختلطة:} الأفضل أن نقسم المهام الكبيرة إلى أوامر متتابعة.
    \item \textbf{قبول أول إجابة:} يمكن دائمًا أن نعدّل الأمر ونعيد المحاولة.
\end{itemize}
\end{intabih}

\begin{arabiyya}[أوامر مفيدة لأستاذ اللغة العربية]
\begin{itemize}
    \item اقترح عشر جمل فيها مفعول لأجله، لتلاميذ السنة الثالثة متوسط.
    \item حوّل هذه الأبيات إلى نثر بسيط، دون تغيير المعنى.
    \item أنشئ خمسة أسئلة فهم حول النص الآتي، مع الإجابات.
\end{itemize}
\end{arabiyya}

\begin{khulasa}
\begin{itemize}
    \item \textbf{الأمر الجيد} يعطي الأداة كل ما تحتاجه لتفهم الطلب.
    \item عناصره الخمسة: \kalima{الدور}، \kalima[akhdar]{المهمة}،
          \kalima[burtuqali]{السياق}، \kalima[ahmar]{القيود}، \kalima[banafsaji]{التنسيق}.
    \item نبني الأمر \textbf{تدريجيًا}، ونحسّنه بعد قراءة الإجابة.
\end{itemize}
\end{khulasa}

\end{document}
